\documentclass[11pt,a4paper]{article}

\usepackage[T1]{fontenc}
\usepackage[utf8]{inputenc}
\usepackage[margin=1in]{geometry}
\usepackage{amsmath,amssymb}
\usepackage{xcolor}
\usepackage{booktabs}
\usepackage{enumitem}
\usepackage{mdframed}
\usepackage{microtype}
\usepackage{parskip}

\definecolor{revisionblue}{RGB}{0,0,255}
\definecolor{commentgray}{RGB}{247,247,247}
\definecolor{commentborder}{RGB}{120,120,120}

\newmdenv[
  backgroundcolor=commentgray,
  linecolor=commentborder,
  linewidth=0.7pt,
  roundcorner=2pt,
  skipabove=8pt,
  skipbelow=8pt,
  innertopmargin=8pt,
  innerbottommargin=8pt,
  innerleftmargin=9pt,
  innerrightmargin=9pt
]{reviewercomment}

\newmdenv[
  linecolor=revisionblue,
  fontcolor=revisionblue,
  linewidth=1pt,
  roundcorner=2pt,
  skipabove=6pt,
  skipbelow=8pt,
  innertopmargin=7pt,
  innerbottommargin=7pt,
  innerleftmargin=9pt,
  innerrightmargin=9pt
]{revisionexcerpt}

\newcommand{\revisionnotice}{\par\noindent\textcolor{revisionblue}{\textbf{Revision highlight in manuscript: blue.}}\par}
\newcommand{\excerptlocation}[1]{\par\noindent\textit{Revised location: #1}\par}

\begin{document}

\begin{center}
{\Large\bfseries Response to Reviewers}\par
\vspace{0.5em}
{\large Manuscript DCHE-D-26-00020R1}\par
\vspace{0.5em}
{\large\bfseries An Interpretable Statistical Surrogate of Activated Sludge Systems}\par
{\large\bfseries that Preserves Mass Conservation and Component Non-Negativity}
\end{center}

\section*{General Manuscript-Wide Revisions}

We thank Reviewer 1 for recognizing the substantial engineering revision and for identifying the remaining points requiring clarification. In response, we have (i) made the symmetry mechanism for the self-quadratic coefficient blocks explicit in both the estimator and manuscript, (ii) corrected the reported-output coefficient extraction so that the interpretability figures represent the raw coupled head $I_{comp}(I_F-\Gamma)^{-1}B$, (iii) reframed every physical guarantee as a property of constrained deployment rather than of the regression head alone, (iv) aligned the Abstract's accuracy description with the quantified moderate predictive tradeoff, and (v) clarified that the accuracy comparison uses a common equality-only affine alignment for comparator outputs, whereas physical-admissibility diagnostics use their direct raw heads. We also regenerated Figure 4 and the supplementary coefficient atlases from the corrected production analysis. The clean and marked manuscripts contain identical scientific content. Blue markup is applied for revisions made directly in response to Reviewer 1's comments/suggestions.

\section*{Reviewer 1}

\subsection*{Overall assessment}

\begin{reviewercomment}
Reviewer 1: The author has made substantial structural revisions to address the concerns raised in the first round. The reformulation from CoBRE, the removal of the bilinear state-input term, the relocation of non-negativity enforcement from a target transformation to a deployment-time linear program, the introduction of a deterministic recursive coupled-QP estimator, the explicit inclusion of HRT as an operating input, the broadened Optuna budget (100 trials per model), and the clarification of the ANN architecture together represent a genuine engineering revision. The mathematical contradiction in the original Equation 5 has been eliminated, and the staged deployment rule (raw $\rightarrow$ affine projection $\rightarrow$ LP) is now stated explicitly. These changes resolve the majority of my earlier concerns.

I have, however, two technical points that still require attention before the manuscript can be recommended for publication, and one framing point.
\end{reviewercomment}

\noindent\textbf{Response.} We sincerely thank the reviewer for recognizing the scope of the first-round revisions and for isolating the remaining technical and framing issues. The new revision addresses both major comments directly and applies the requested accuracy framing consistently across the submission.

\subsection*{Major comment 1}

\begin{reviewercomment}
\textbf{Major comment 1: The symmetry of $\Theta_{cc}$ and $\Theta_{uu}$ is asserted but not demonstrated}

In response to my original R1.33, the authors state that the influent-influent block $\Theta_{cc}$ and the operational-operational block $\Theta_{uu}$ are symmetric ``by construction'' and that no $(\Theta+\Theta^{T})/2$ averaging is therefore required before the heatmaps in Figure 4 and the supplementary atlases can be interpreted cell-by-cell. The claim is consequential: Section 4.4 reads off specific numerical entries ($S_O$-$S_O$ = $-2.3899\times10^{-2}$, $S_O$-$S_{N_2}$ = $-2.1845\times10^{-2}$, $S_O$-$S_{NO_2}$ = $-1.4691\times10^{-2}$) and assigns physical meaning to them.

The text supporting this claim, however, is not adequate as currently written. The Nomenclature lists $\Theta_{cc}\in\mathbb{R}^{F\times F^2}$ and $\Theta_{uu}\in\mathbb{R}^{F\times M_{op}^2}$, i.e., the general non-symmetric parameterization of a second-order tensor vectorized through the Kronecker product. Because $c_{in}\otimes c_{in}$ contains the products $c_i c_j$ and $c_j c_i$ as separate entries that are numerically equal, the effect on each output target is the sum $\theta_{k,ij}+\theta_{k,ji}$. Without an explicit constraint imposed during fitting, that decomposition is not unique: infinitely many $\theta_{k,ij},\theta_{k,ji}$ pairs reproduce the same prediction, so the individual cells of the heatmap have no well-defined physical meaning even though the model fit is perfectly valid.

The response argues that the partitioned feature map and the corresponding driver coefficients are ``organized so that the fitted $\Theta$ blocks satisfy the commutativity-induced symmetry exactly,'' but the manuscript does not actually describe the mechanism that produces this symmetry. The reader cannot determine whether the authors (i) parameterize only the upper-triangular indices and recover the lower-triangular entries by mirroring, (ii) impose $\theta_{k,ij}=\theta_{k,ji}$ as an equality constraint inside the QP, or (iii) project onto the symmetric subspace after each block-coordinate update. Any one of these would resolve the issue; none is currently in the text.

I ask the authors to add an explicit technical sentence to Section 2.3 (or to the corresponding subsection of 2.6) stating how symmetry is enforced, with a one-line justification why this preserves the optimum of the coupled QP. If the code does not in fact enforce symmetry, then either a symmetrization layer must be added or the symmetric average $(\Theta+\Theta^{T})/2$ must be plotted before any cell-level physical interpretation is made in Section 4.4.
\end{reviewercomment}

\noindent\textbf{Response.} We are grateful to the reviewer for identifying that the previous explanation did not establish coefficient identifiability. We implemented option (iii): after the initial ridge solve and after every recursive $B$ update, every output-specific $\Theta_{uu}$ and $\Theta_{cc}$ slice is explicitly projected onto the symmetric subspace.

Section 2.3 now defines the square output-specific reshapes, states the projection, and supplies the requested optimum-preservation proof. Section 2.6 and Pseudocode 1 now state that symmetrization follows every $B$ update, and the Nomenclature identifies both blocks as symmetrized. The Figure 4 caption now states that off-diagonal cells are mirrored and that each displayed off-diagonal cell is half the combined unordered cross-term coefficient. Separately, the coefficient-retention table caption and the paragraph following that table state that retention counts use only unique upper-triangular terms; the heatmap itself continues to display the full mirrored matrix.

We revised the Python workflow implementing this. The revised workflow first constructs the globally defined raw-head reported-output map $I_{comp}\widehat{R}^{-1}\widehat{B}$ and then extracts COD, TN, TP, and TSS blocks. Figure 4, the coefficient-retention table, and Supplementary Figures S1--S4 were regenerated from this corrected production mapping. The corrected COD raw-head values for the three pairs cited by the reviewer are $S_O$--$S_O=-5.8097$, $S_O$--$S_{N2}=-1.7411$, and $S_O$--$S_{NO2}=-9.4127\times10^{-2}$; they supersede the former values. The manuscript explicitly limits these coefficient interpretations to the pre-correction raw head because the final affine/LP map is piecewise and has no single global coefficient matrix.

\revisionnotice

\excerptlocation{Nomenclature}
\begin{revisionexcerpt}
Influent--influent driver block ($\Theta_{cc} \in \mathbb{R}^{F \times F^2}$); each output-specific $F \times F$ reshape is symmetrized

Operating--operating driver block ($\Theta_{uu} \in \mathbb{R}^{F \times M_{op}^2}$); each output-specific $M_{op} \times M_{op}$ reshape is symmetrized
\end{revisionexcerpt}

\excerptlocation{Section 2.3, ``Coupled Second-Order Driver and Prediction-Space Training''}
\begin{revisionexcerpt}
For output component $k$, let $Q_{uu}^{(k)} \in \mathbb{R}^{M_{op} \times M_{op}}$ and $Q_{cc}^{(k)} \in \mathbb{R}^{F \times F}$ denote the square reshapes of the corresponding rows of $\Theta_{uu}$ and $\Theta_{cc}$. The full ordered Kronecker map retains both $x_i x_j$ and $x_j x_i$, so after every $B$ update the implementation projects each self-quadratic slice onto the symmetric subspace,
\begin{equation*}
Q \leftarrow \operatorname{sym}(Q) = \frac{Q + Q^T}{2},
\qquad Q \in \left\{Q_{uu}^{(k)},Q_{cc}^{(k)}\right\}_{k=1}^{F}.
\end{equation*}
This projection preserves the coupled-QP optimum. For any vector $x$, $x^TQx=x^T\operatorname{sym}(Q)x$, whereas
\begin{equation*}
\|Q\|_F^2 = \|\operatorname{sym}(Q)\|_F^2 + \|\operatorname{skew}(Q)\|_F^2.
\end{equation*}
Consequently, removing the skew-symmetric part leaves every fitted value unchanged and cannot increase the ridge-regularized objective. Because the reported estimator uses $\lambda_B>0$, the unique ridge solution has zero skew part, so the explicit projection makes the commutativity-induced equality $\theta_{k,ij}=\theta_{k,ji}$ hold to machine precision.
\end{revisionexcerpt}

\excerptlocation{Section 2.6, ``Deterministic Recursive Coupled-QP Estimation and Modeling Boundaries''}
\begin{revisionexcerpt}
For fixed $(\Gamma, \widehat{C})$, the $B$ subproblem is a ridge-regularized linear least-squares solve followed by the prediction-preserving symmetric projection stated in Section~2.3. For fixed $(B, \widehat{C})$, the $\Gamma$ subproblem is a convex box-constrained quadratic program. For fixed $(B, \Gamma)$, the $\widehat{C}$ subproblem is a convex non-negative quadratic program in which invariant mismatch remains a soft quadratic penalty.
\end{revisionexcerpt}

\excerptlocation{Pseudocode 1}
\begin{revisionexcerpt}
compute $B^{(0,s)}$ from the corresponding ridge solve and project every $Q_{uu}^{(k)}$ and $Q_{cc}^{(k)}$ onto the symmetric subspace

update $B^{(t+1,s)}$ via Eq.~(14) and project every $Q_{uu}^{(k)}$ and $Q_{cc}^{(k)}$ onto the symmetric subspace
\end{revisionexcerpt}

\excerptlocation{Section 4.4, ``Interpretable Raw-Head Structure of the Learned ICSOR Model''}
\begin{revisionexcerpt}
Coefficient interpretation is performed on the globally defined pre-correction head rather than on the piecewise deployment correction. Since $c_{raw}=\widehat{R}^{-1}\widehat{B}\phi$, its reported-composite counterpart is
\begin{equation*}
y_{raw}=I_{comp}c_{raw}=B_{raw,ext}\phi,
\qquad
B_{raw,ext}=I_{comp}\widehat{R}^{-1}\widehat{B}.
\end{equation*}
Figure~4 presents the COD row's $\Theta_{cc}$ block from $B_{raw,ext}$. The COD, TN, TP, and TSS raw-head atlases are provided as Supplementary Figures S1--S4. Rows of $\widehat{B}$ are ASM component-driver rows and therefore cannot be indexed directly by reported-composite labels. The final affine/LP map is active-set dependent and has no single global coefficient matrix; Figure~4 must not be read as a derivative of the final deployed output.
\end{revisionexcerpt}

\excerptlocation{Figure 4 caption}
\begin{revisionexcerpt}
COD raw-head $\Theta_{cc}$ heatmap from $B_{raw,ext}=I_{comp}\widehat{R}^{-1}\widehat{B}$. Rows and columns enumerate influent ASM components. The matrix is symmetric to machine precision; off-diagonal cells are mirrored and each displayed cell is half the coefficient of the corresponding unordered cross term. The full raw-head atlases are reported in Supplementary Figures S1--S4.
\end{revisionexcerpt}

\excerptlocation{Section 4.4 coefficient interpretation}
\begin{revisionexcerpt}
The COD raw-head surface is exactly symmetric at stored precision. Its largest-magnitude entry is the $S_O$--$S_O$ self-curvature term ($-5.8097$), followed by the mirrored $S_O$--$S_{N2}$ pair ($-1.7411$). Other prominent terms link $S_O$ with $X_{MeP}$ ($2.3190\times10^{-1}$), $X_{AOB}$ ($-1.3264\times10^{-1}$), $S_{PO4}$ ($-1.3055\times10^{-1}$), $X_{NOB}$ ($-1.0182\times10^{-1}$), and $S_{NO2}$ ($-9.4127\times10^{-2}$). Under the ordered Kronecker convention, a displayed off-diagonal value $q_{ij}=q_{ji}$ contributes $(q_{ij}+q_{ji})x_ix_j=2q_{ij}x_ix_j$ to the polynomial, so cell values and combined unordered cross-term coefficients are not interchangeable.
\end{revisionexcerpt}

\excerptlocation{Coefficient-retention table caption and entries}
\begin{revisionexcerpt}
Retained coefficients in the COD raw-head structure. The reported-output blocks use $10^{-4}$ of the largest absolute COD raw-head coefficient ($4.7332\times10^{-3}$); the shared $\Gamma$ block uses $10^{-4}$ of its own largest off-diagonal magnitude ($3.6803\times10^{-5}$). Symmetric quadratic totals count only unique upper-triangular terms.

\begin{center}
\scriptsize
\begin{tabular}{lrrr}
\toprule
Coefficient block & Total coeffs. & Retained coeffs. & \% Retained \\
\midrule
$b$ & 1 & 1 & 100.0 \\
$W_u$ & 2 & 2 & 100.0 \\
$W_{in}$ & 20 & 19 & 95.0 \\
$\Theta_{uu}$ & 3 & 3 & 100.0 \\
$\Theta_{uc}$ & 40 & 29 & 72.5 \\
$\Theta_{cc}$ & 210 & 44 & 21.0 \\
$\Gamma$ & 380 & 371 & 97.6 \\
Overall summary & 656 & 469 & 71.5 \\
\bottomrule
\end{tabular}
\end{center}
\end{revisionexcerpt}

\excerptlocation{Paragraph following the coefficient-retention table}
\begin{revisionexcerpt}
Table~12 counts independently interpretable terms rather than duplicated display cells. Thus $\Theta_{uu}$ has $2(2+1)/2=3$ unique terms and $\Theta_{cc}$ has $20(20+1)/2=210$, although their displayed heatmaps contain 4 and 400 cells, respectively. The reporting threshold is post-training only and does not affect model selection, fitting, or deployment. Nineteen of 20 $W_{in}$ entries, all three unique $\Theta_{uu}$ terms, 29 of 40 $\Theta_{uc}$ terms, and 44 of 210 unique $\Theta_{cc}$ terms exceed the COD raw-head threshold; 371 of 380 off-diagonal $\Gamma$ entries exceed the separately scaled coupling threshold. Overall, 469 of 656 unique candidates (71.5\%) are retained. The result indicates selective influent curvature within a comparatively dense coupling structure, not a uniformly sparse model.
\end{revisionexcerpt}

\subsection*{Major comment 2}

\begin{reviewercomment}
\textbf{Major comment 2}

Section 4.3 reports honestly that none of the raw coupled predictions are both invariant-consistent and non-negative (raw conservation compliance 0.0\%, raw non-negativity 21.7\%). The 0\% violation rate is therefore not a property of the learned regression, the regression head itself violates conservation on 100\% of test samples, but of the deployment correction: the affine projection resolves 21.8\% of cases and the LP handles the remaining 78.2\%. The invariant penalty $\lambda_{inv}$ in Equation 7 is a soft regularizer during training; the hard guarantee comes from Equation 12.

The Abstract, Graphical Abstract, and Conclusion currently project the 0\% outcome as if it emerged from the regression. I recommend a brief reframing across these sections so that ICSOR is described as a constrained-deployment surrogate whose regression head is structured to remain compatible with hard conservation correction, together with a single sentence in Section 2.4 stating that the hard guarantee is produced by the deployment correction rather than by the training objective.
\end{reviewercomment}

\noindent\textbf{Response.} We thank the reviewer for this precise distinction and agree. The Abstract, Highlights, Graphical Abstract, Introduction, Sections 2.3, 2.4, and 2.6, Sections 4.1, 4.3, and 4.4, and the Conclusion now consistently describe ICSOR as an interpretable constrained-deployment surrogate. The regression head is described as invariant-aware but not inherently feasible. Section 2.4 now states explicitly that $\lambda_{inv}$ is a soft regularizer and that the hard guarantee is supplied by the checked affine/LP deployment rule. Section 4.3 now states directly that the zero final violation rates belong to the ICSOR deployment map, not to its regression head.

The revised Graphical Abstract uses the exact labels ``Invariant-aware regression head,'' ``soft training penalties,'' ``Hard deployment correction:,'' and ``Raw $\rightarrow$ Affine $\rightarrow$ LP.'' Its outcome panel is headed ``Final deployed outputs:'' and reports ``0\% conservation violations'' and ``0\% non-negativity violations,''.

\revisionnotice

\excerptlocation{Abstract}
\begin{revisionexcerpt}
Activated sludge models are essential for wastewater-treatment design and operation, but their detailed simulations can be too computationally demanding for repeated analysis. Machine-learning surrogates are faster, yet accurate predictions may still violate mass conservation or produce negative component concentrations. This study introduces Invariant-Constrained Second-Order Regression (ICSOR), an interpretable surrogate that combines a transparent, invariant-aware regression head with a constrained deployment procedure. The regression head predicts the effluent component state from influent conditions and operating inputs. At deployment, each prediction is checked, projected onto the mass-conservation constraints, and, when necessary, adjusted by a linear program to enforce both conservation and non-negativity. ICSOR was evaluated on 10,000 steady-state samples from an ASM2d-TSN continuous stirred tank reactor and compared with XGBoost, LightGBM, CatBoost, AdaBoost, Random Forest, support vector regression, $k$-nearest neighbors, partial least squares, and a multilayer perceptron. The deployed ICSOR predictions had zero conservation and non-negativity violations: the affine projection resolved 21.8\% of test cases, and the linear program resolved the remaining 78.2\%. For context, none of the raw regression-head predictions satisfied conservation, although 21.7\% were non-negative, showing that the hard guarantee came from constrained deployment rather than training alone.
\end{revisionexcerpt}

\excerptlocation{Highlights}
\begin{revisionexcerpt}
ICSOR combines interpretable regression with constrained deployment.

Deployed predictions satisfy conservation and non-negativity exactly.

Physical guarantees involve a moderate fixed-split accuracy tradeoff.
\end{revisionexcerpt}

\excerptlocation{Graphical Abstract labels}
\begin{revisionexcerpt}
Invariant-aware regression head

soft training penalties

Hard deployment correction:

Raw $\rightarrow$ Affine $\rightarrow$ LP

Final deployed outputs:

0\% conservation violations

0\% non-negativity violations

\end{revisionexcerpt}

\excerptlocation{Introduction}
\begin{revisionexcerpt}
This study addresses that need through Invariant-Constrained Second-Order Regression (ICSOR), an interpretable surrogate for steady-state activated sludge systems. ICSOR is designed for a practical task: using influent conditions and operating settings to provide a fast estimate of the effluent state that engineers can inspect and use with confidence. Unlike an unconstrained predictor, ICSOR checks the predicted component state before returning it and ensures that the final result conserves the modeled material inventories and contains no negative concentrations.
\end{revisionexcerpt}

\excerptlocation{Section 2.3, training-objective interpretation}
\begin{revisionexcerpt}
Each term has a distinct role. The first term fits effluent ASM components directly. The second softly penalizes invariant mismatch in component space; it is not a hard equality constraint. The third encourages consistency between the fitted states and the coupled driver relation. The ridge penalties stabilize the driver and coupling estimates, and the positive $B$ ridge penalty also selects the symmetric minimum-norm representative of each duplicated ordered quadratic pair.
\end{revisionexcerpt}

\excerptlocation{Section 2.4, ``Deployment-Time Inference with Exact Conservation and Non-Negativity''}
\begin{revisionexcerpt}
Equation~(12) is the final correction branch of the checked deployment rule.

The hard feasibility guarantee is produced by this checked deployment rule, not by the training objective. In particular, the $\lambda_{inv}$ term in Eq.~(7) is a soft regularizer and does not guarantee that $c_{raw}$ satisfies $Ac_{raw}=Ac_{in,*}$. A raw or affine state is returned only after it passes the applicable feasibility checks; otherwise Eq.~(12) imposes $Ac=Ac_{in,*}$ and $c\geq 0$ jointly. Because the physically admissible $c_{in,*}$ is itself feasible under the adopted component basis, the LP feasible set is nonempty.
\end{revisionexcerpt}

\excerptlocation{Section 2.6, estimator description}
\begin{revisionexcerpt}
For fixed $(B, \Gamma)$, the $\widehat{C}$ subproblem is a convex non-negative quadratic program in which invariant mismatch remains a soft quadratic penalty.
\end{revisionexcerpt}

\excerptlocation{Section 3.3, benchmark protocol}
\begin{revisionexcerpt}
ICSOR differs in two distinct respects that should not be conflated: its regression head is trained with soft invariant and coupled-system penalties and a non-negative auxiliary fitted state, while exact conservation and non-negativity are imposed only on the final state returned by its deployment correction.
\end{revisionexcerpt}

\excerptlocation{Section 4.1, ``Overall Benchmark Performance on the Fixed Train-Test Split''}
\begin{revisionexcerpt}
In terms of predictive accuracy on the fixed split, the MLP attains the lowest error among the retained baselines and LightGBM attains the lowest among the tree-based methods. ICSOR does not attain the lowest aggregate RMSE: its final deployed outputs give 5.98, which is 36.5\% higher than the MLP value of 4.38 and 12.8\% higher than the LightGBM value of 5.30. ICSOR remains within 0.01 of SVR and below XGBoost and CatBoost. These results define a moderate accuracy--feasibility tradeoff.
\end{revisionexcerpt}

\excerptlocation{Section 4.3, ``Physical Admissibility and Deployment Behavior''}
\begin{revisionexcerpt}
The staged deployment logic can also be examined at the level of individual correction stages. None of the raw coupled test predictions was simultaneously invariant-consistent and non-negative: raw feasibility was 0.0\%, raw conservation compliance was 0.0\%, and 21.7\% of raw states was non-negative. The closed-form affine projection restored invariants for all samples and was itself sufficient for 21.8\% of the test set (436 of 2000 samples), resolving those cases without invoking a numerical optimizer. The remaining 78.2\% (1564 samples) required the final LP correction, which converged in a mean of 23.5 HiGHS iterations with a maximum of 40. The final returned prediction was feasible for 100\% of the test cases. Thus, the zero final violation rates are properties of the constrained deployment map, not of the learned regression head in isolation.
\end{revisionexcerpt}

\excerptlocation{Section 4.4, raw-head framing}
\begin{revisionexcerpt}
This coefficient-level view differs from PLS loadings and local ANN sensitivity analysis because the ICSOR raw head exposes named global blocks for operating effects, loading effects, within-subspace curvature, and cross-component coupling. The values are shaped by invariant-aware soft training and the coupled non-negative fitted-state variable, but raw-head feasibility is not guaranteed. Exact physical admissibility is imposed afterward by constrained deployment. The coefficient blocks therefore support transparent structural inspection without implying that the raw head is itself feasible.
\end{revisionexcerpt}

\excerptlocation{Conclusion}
\begin{revisionexcerpt}
This study introduced Invariant-Constrained Second-Order Regression (ICSOR), an interpretable surrogate for steady-state activated sludge systems. ICSOR separates prediction from physical enforcement. Its invariant-aware regression head provides a transparent component-level prediction, and a hierarchical deployment procedure ensures that the returned state satisfies exact stoichiometric conservation and componentwise non-negativity. The procedure first checks the raw prediction, then applies a closed-form affine projection, and finally uses a linear program when non-negativity is still violated. The affine stage was sufficient for 21.8\% of test cases, while the linear program handled the remaining 78.2\%. Consequently, every deployed ICSOR prediction satisfied both physical requirements. By comparison, every unconstrained baseline violated conservation on all test samples, and several also predicted negative component concentrations.
\end{revisionexcerpt}

\subsection*{Minor comment}

\begin{reviewercomment}
\textbf{Minor comment}

ICSOR's aggregate test RMSE of 5.98 is 36\% higher than the MLP's 4.38 and 13\% higher than LightGBM's 5.30. Calling this ``competitive aggregate accuracy'' in the Abstract is generous; the Conclusion describes the same result more carefully as a ``moderate predictive tradeoff.'' I recommend aligning the Abstract with the more careful framing in the Conclusion so that the reader is not surprised by the accuracy gap when reaching Section 4.1.
\end{reviewercomment}

\noindent\textbf{Response.} We appreciate the reviewer's recommendation and have adopted the more cautious framing. The Abstract no longer claims ``competitive aggregate accuracy.'' It now reports the three RMSE values directly and describes the result as a moderate predictive tradeoff. Section 4.1 quantifies the gaps as 36.5\% relative to the MLP and 12.8\% relative to LightGBM under the stated reporting protocol, in which comparator values include the common equality-only affine alignment but no LP non-negativity correction; the Conclusion uses the same moderate-tradeoff description.

\revisionnotice

\excerptlocation{Abstract}
\begin{revisionexcerpt}
This physical reliability involved a moderate accuracy tradeoff: aggregate test RMSE was 5.98, compared with 4.38 for the multilayer perceptron and 5.30 for LightGBM.
\end{revisionexcerpt}

\excerptlocation{Section 4.1, ``Overall Benchmark Performance on the Fixed Train-Test Split''}
\begin{revisionexcerpt}
ICSOR does not attain the lowest aggregate RMSE: its final deployed outputs give 5.98, which is 36.5\% higher than the MLP value of 4.38 and 12.8\% higher than the LightGBM value of 5.30. ICSOR remains within 0.01 of SVR and below XGBoost and CatBoost. These results define a moderate accuracy--feasibility tradeoff. Under the stated reporting protocol, the comparator values include the common equality-only affine alignment but no LP non-negativity correction, whereas ICSOR is evaluated after its complete checked deployment rule.
\end{revisionexcerpt}

\excerptlocation{Conclusion}
\begin{revisionexcerpt}
These physical guarantees came with a moderate accuracy tradeoff. On the fixed 80--20 split, ICSOR produced an aggregate test RMSE of 5.98, compared with 4.38 for the MLP and 5.30 for LightGBM. These differences correspond to increases of 36.5\% and 12.8\%, respectively.
\end{revisionexcerpt}

\subsection*{Acceptance condition}

\begin{reviewercomment}
Conditional on Major comments 1 and 2 being addressed, the manuscript will be in shape for acceptance.
\end{reviewercomment}

\noindent\textbf{Response.} We sincerely thank the reviewer for this assessment. We have addressed both major comments in the estimator, analysis workflow, manuscript, figures, supplementary material, and marked revision, and we hope the revised submission now meets the standard for acceptance.

\end{document}
